% Experiments --- setup, datasets, metrics, harness. Numbers live in
% Results; this section fixes the protocol. The LLM-benchmark table is a
% honest placeholder pending US-069 (H100 eval), per the data-real rule.
\section{Experimental Setup}
\label{sec:experiments}

\subsection{Datasets}
\paragraph{PASTIS-R}~\cite{garnot2021pastis} is the primary benchmark:
2{,}433 multi-temporal Sentinel-2 patches ($128\times128$ px), 18 crop
classes, metropolitan France, with five spatially disjoint folds.
Unless stated otherwise, segmentation models are evaluated on a
held-out fold and ensemble combiners on out-of-fold predictions.

\paragraph{Sen4AgriNet (Catalonia)}~\cite{sykas2022sen4agrinet}
(CC-BY-SA-4.0) is used for the France$\rightarrow$Catalonia transfer
study. A Catalonia subset is versioned with DVC
(\texttt{data/sen4agrinet.dvc}).

\paragraph{EuroCropsML}~\cite{reuss2025eurocropsml} (CC-BY-SA-4.0)
supports the few-shot $k$-shot curve; the raw subset is versioned with
DVC (\texttt{data/transfer/eurocropsml.dvc}).

\paragraph{AgroMind}~\cite{ruan2025agromind} is used only to evaluate
the conversational layer. It has no training split (about 28{,}482 QA
pairs); fine-tuning on it would induce leakage, so it is strictly an
evaluation set.

\subsection{Metrics}
Segmentation is reported with mean intersection over union (mIoU) and
macro-averaged F1 at the parcel level; we also report pixel accuracy
where relevant. Macro averaging is the primary objective because it
penalises errors on rare and common crops equally, matching the
business objective. Ensemble combiners are compared by out-of-fold
macro F1 and accuracy.

\subsection{Evaluation harness and reproducibility}
All segmentation runs are scored by a shared evaluation harness
(\texttt{ml/eval/}) that enforces spatially disjoint folds and aborts on
geographic overlap between training and evaluation sets. A label-remap
module (\texttt{ml/eval/class\_remap.py}) projects predictions onto the
harmonised macro label space for cross-region comparison. Every figure
reported in Section~\ref{sec:results} is anchored to a versioned
artefact (JSON/CSV/registry) in the repository, recorded in
\texttt{\% src:} comments in the source; datasets and weights are
tracked with DVC and runs with MLflow.

\subsection{Conversational-layer benchmark protocol}
\label{sec:experiments-llm}
The conversational layer is evaluated on AgroMind with an
LLM-as-judge protocol comparing the frozen Gemini~2.5~Pro reasoner
against the on-premises Qwen multimodal reasoner served via vLLM. The
quantitative benchmark table is produced by a separate evaluation run
on the H100 GPU (US-069) and is intentionally left as a structured
placeholder here; numbers are not fabricated.
Table~\ref{tab:llm-bench} fixes the column structure.

\begin{table}[htbp]
\caption{Conversational-layer benchmark (structure only; numbers
pending the H100 evaluation, US-069).}
\label{tab:llm-bench}
\centering
\begin{tabular}{lccc}
\toprule
\textbf{Reasoner (frozen)} & \textbf{Accuracy} & \textbf{Grounding}
  & \textbf{LLM-judge} \\
\midrule
Gemini~2.5~Pro (cloud)      & --- & --- & --- \\
Qwen3.5-35B-A3B (on-prem)   & --- & --- & --- \\
\bottomrule
\end{tabular}
% src: pending US-069 (H100 eval) -- do not fill with synthetic numbers.
\end{table}
